\chapter*{Scope, Primitives, and Claim Semantics}
\addcontentsline{toc}{chapter}{Scope, Primitives, and Claim Semantics}

\section*{What is and is not being derived}

Throughout this monograph, the word \emph{derived} has a conditional mathematical meaning. Temporal structures are derived from a declared relational model and its typed primitives; they are not claimed to follow uniquely from Shannon information alone, nor is the existence of physical time inferred from an information measure without additional assumptions.

The temporal branch begins from the following admitted objects:
\begin{enumerate}
\item a set of relational states and admitted directed transitions;
\item a normalized probability vector \(p(s)\) on each state, with Shannon entropy \(H_S\);
\item a normalized complex representative \(\psi(s)\) on sectors where phase transport is defined;
\item positive admitted transition probabilities or rates, together with declared relational mobility fields;
\item an ordering parameter \(\lambda\), which is \emph{not} metric clock time;
\item the inherited normalization \(\kappa=\ln 2/(24\pi)\).
\end{enumerate}

From those inputs the programme asks which further structures follow by definition, algebraic identity, dynamical closure, or experimentally testable bridge. The positive elapsed measure, relative lapse, memory carriers, retrodiction maps, and later spacetime interface are therefore downstream constructions whose validity is conditional on their stated domains and dependencies.

\section*{Standard ingredients and the IDT-specific composition}

Several ingredients used in the construction are standard and are not presented as new results. These include Shannon entropy \cite{Shannon1948}, Kullback--Leibler divergence \cite{KullbackLeibler1951}, Pancharatnam/Berry geometric phase \cite{Pancharatnam1956,Berry1984}, quantum-state information geometry \cite{ProvostVallee1980}, Onsager reciprocity and nonequilibrium network structure \cite{Onsager1931,Schnakenberg1976}, and entropy-gradient formulations for reversible finite Markov chains \cite{Maas2011}. The relativistic closure layer likewise uses standard Noether and Einstein structures \cite{Noether1971,Einstein1916}.

The programme-specific content is the typed way these ingredients are composed and constrained: the separation of exact state entropy from non-exact path affinity, the positive-activity temporal measure, the clock-ratio and memory interfaces, the constructive retrodiction carrier, the evidence-class firewall between retrodiction and retrocausal claims, and the explicit temporal-to-spacetime closure gates. A standard identity remains standard even when it is used as an important recovery target inside this architecture.

\section*{Evidence classes}

Every substantive statement belongs to one of five evidence classes:
\[
\boxed{
\begin{array}{c}
\text{definition / admitted primitive}\\
\downarrow\\
\text{proved identity on a declared domain}\\
\downarrow\\
\text{reference-model or software result}\\
\downarrow\\
\text{calibrated physical correspondence}\\
\downarrow\\
\text{replicated measured physical result}.
\end{array}}
\]
A result does not move upward merely because more unit tests pass. Hosted tests establish implementation consistency and regression protection; they are not independent physical observations.

\section*{Theory map}

The monograph follows four layers:
\[
\boxed{
\begin{aligned}
\textbf{I. Foundations:}\quad&
\text{relational state}
\to
\text{information + phase}
\to
\text{positive temporal activity},
\\[1mm]
\textbf{II. Temporal dynamics:}\quad&
\text{elapsed measure}
\to
\text{NOW/bifurcation}
\to
\text{transport + memory},
\\[1mm]
\textbf{III. Inference and experiment:}\quad&
\text{retrodiction}
\to
\text{withheld-variable tests}
\to
\text{preregistered physical tests},
\\[1mm]
\textbf{IV. Spacetime closure:}\quad&
\text{temporal coframe}
\oplus
\text{independent spatial/source branch}
\to
\text{Einstein interface}.
\end{aligned}}
\]

This map also fixes the novelty boundary. Standard-physics relations used at Layer IV are recovery benchmarks unless the temporal branch supplies an independently determined quantity that can disagree with the standard prediction before the measurement is inspected.
